% Body of Section 2; \section{Methods} is set in main-sc.tex.

\subsection{Process design}\label{sec:methods-process}

\subsubsection{Main process unit}\label{sec:methods-main}

\begin{figure}[htbp]
    \centering
    \includegraphics[width=\textwidth]{BiomassLCA-BioPyrolysis (3).pdf}
    \caption{Biomass pyrolysis-to-SAF production process flow diagram. The diagram shows the major process units, biomass and hydrogen feed, intermediate streams, and product outputs in the analysis.}
    \label{fig:pfd-main}
\end{figure}

The process is based on the biomass fast-pyrolysis-to-fuels flowsheet developed by NREL/PNNL \citep{jones2013process}. The flowsheet comprises eight process areas: feed handling, pyrolysis, hydrotreating, product finishing, hydrogen recovery, hydrocracking, off-gas combustion, and the steam and balance-of-plant island (Fig.~\ref{fig:pfd-main}). Detailed description is available in SI~S1. 

In this study, the process design is adjusted from the early study of \citep{jones2013process} in four distinctive points. 
First, product slate of the process is shifted to maximize SAF production for the South Korean market, rather than the gasoline and diesel blendstock production targeted in the original design. Gasoline- and diesel-range streams are retained as value-generating co-products.
Second, surplus electricity is generated as a co-product via a gas-turbine power island, which is newly introduced in this study. 
Next, the original grey-hydrogen supply can be partially or fully replaced with hydrogen produced by green-ammonia cracking, as described in Section~\ref{sec:methods-cracker}.
Lastly, an amine-based CCS unit is added to capture recoverable CO$_2$ streams leaving combustion chambers or separators (Section~\ref{sec:methods-ccs}).

At the base scale, the plant processes 0.343~MMt/yr of woody biomass and produces 0.041~MMt/yr of SAF. The bio-oil sent to hydrotreating amounts to 0.244~MMt/yr, equivalent to approximately 71~wt\% of the biomass feed. Total fresh make-up hydrogen consumption is 0.024~MMt/yr, comprising 0.021~MMt/yr for hydrotreating and 0.003~MMt/yr for hydrocracking. Hydrotreating employs ruthenium catalysts in a series of two reactors at a liquid hourly space velocity (LHSV) of 0.5, followed by the last reactor with molybdenum catalyst at an LHSV of 0.22. Catalyst costs are estimated separately from reactor-vessel costs.

The Aspen Plus flow-sheet is costed using Aspen Process Economic Analyzer (APEA). The resulting \textit{Main} block, which represents whole-flow sheet capital and operating expenditures (CAPEX and OPEX) obtained from the APEA, is combined with five separately costed sub-units: CCS, pressure swing adsorption (PSA) hydrogen recovery, the off-gas reformer, hydrotreating catalysts, and the hydrocracker. Cost interpolation for these sub-units is detailed in SI~S3.

Plant capacity is scaled from the base-case SAF output of 0.041 MMt/yr to 1.23 MMt/yr, representing a 30-fold difference between the smallest and largest plant scales. Rather than applying a single cost-scaling exponent to the entire plant, each major unit is assigned an explicit single-train capacity limit. Additional parallel trains are introduced when a unit reaches its limit. The capacity limits and corresponding numbers of trains at each scale are reported in SI~S4. Because additional trains replicate existing units rather than enlarge individual units, they do not benefit from further single-unit economies of scale. This train-level configuration is implemented explicitly in the assessment at each plant scale.

All costs are reported in first-quarter 2022 USD. Costs from the original study of \citep{jones2013process} are escalated using a cost-index ratio of 3230/2250. The financial assumptions are 8,000 operating hours per year, a 30-year plant life, a 10\% internal rate of return (IRR), and a 35\% tax rate. The complete cost basis and financial assumptions are provided in SI~S2.

\subsubsection{Ammonia cracker}\label{sec:methods-cracker}

Green ammonia is used as a hydrogen carrier, with ammonia imported and catalytically cracked at the site to supply make-up hydrogen. The cracker capacity is determined by the plant's make-up hydrogen demand, with a single-train capacity limit of 403~GWh/yr. The reaction network, gas hourly space velocity (GHSV), and kinetics follow the source used to develop the Aspen model \citep{cho2024ammonia}. Ammonia conversion and hydrogen recovery also follow the parameters set in the original source. 

The recovery process was modified to employ a sulfuric acid scrubbing system to capture converted ammonia. It employs a multi-stage sulfuric acid scrubber to convert left ammonia into a liquid fertilizer \citep{manuzon2007prototype}. At the end, the scrubbed gas proceeds to nitrogen gas separation, yielding high-purity hydrogen (99.9\% H$_2$). The produced nitrogen fertilizer is credited as an ammonia-side co-product. Further details of the cracker process are provided in SI~S5.

\subsubsection{Carbon capture and storage}\label{sec:methods-ccs}

Post-combustion CCS is modeled using a monoethanolamine (MEA)-based retrofit. Its cost is interpolated from the power-plant amine-capture reference case reported by \citet{rao2002}, which processes 3.13~MMt/yr of CO$_2$, has a total CAPEX of \$297.9~M, requires 7,772~GWh/yr of regeneration heat, and consumes auxiliary electricity equivalent to 14\% of the heat duty. All costs are converted to refer to first-quarter 2022 USD.

The total captured mass comprises the exhaust-gas-recirculation (EGR) CO$_2$ stream and a specified fraction of the recoverable process CO$_2$ stream. The EGR CO$_2$ stream is completely captured. Capture-unit CAPEX scales with captured mass according to the six-tenths rule, whereas OPEX scales with an exponent of 0.7. Utility requirements scale linearly with captured mass. Capture OPEX includes make-up fuel for regeneration heat not supplied by on-site heat, auxiliary electricity, and optional CO$_2$ transport and storage (T\&S), the details for which is described in section~\ref{sec:methods-ccscost}). The capture fraction is varied from 0 to 1.0. The complete CCS cost derivation is provided in SI~S6.

\subsection{Techno-economic analysis}\label{sec:methods-tea}

The economic performance metric is the minimum jet-fuel selling price (MJSP,  \$/L SAF). MJSP is determined by discounted cash-flow analysis as the SAF selling price required to achieve a zero net present value given the economic parameters, assuming straight-line depreciation and zero salvage value.
Revenue from naphtha, diesel, exported electricity, and ammonia-side fertilizer is credited against the required revenue. Electricity and natural gas prices are set at \$65.0/MWh and \$49.1/MWh, respectively.

\subsubsection{Feedstock price model}\label{sec:methods-feed}

Because the domestic forest-residue supply is finite, the delivered biomass price is modeled as a function of plant scale rather than held constant. Larger plants are assumed to require progressively more expensive feedstock as regional supply is depleted and additional sources are accessed. 
% The regional supply context is described in SI~S8. 

% The scarcity of FR biomass feedstock was reflected in the scale-up analysis as a feedstock price via two main approaches. 
This study considers the full spectrum of woody biomass as viable feedstock, capitalizing on the proven versatility of pyrolysis to process both softwood and hardwood species. Literature consistently demonstrates successful bio-oil conversion across a diverse range of softwoods, such as pine and conifers \cite{tran2021fast, jarvis2014characterization}, as well as hardwoods, including black locust and birch \cite{straker2015black, li2014autothermal}. While comparative research indicates that compositional factors—specifically differences in lignin and ash content—may influence the precise yield profile \citep{garcia2007vacuum, azeez2010fast}, the fundamental feasibility of thermochemical conversion remains robust across species. Consequently, this process is designed to be feedstock-flexible, accommodating any woody biomass regardless of its classification.

The economic viability of this feedstock-flexible approach is further supported by the immense volume of available biomass. The total potential supply is substantial, comprising approximately 706~MMt of logs and 412~MMt of underutilized residues. However, these resources are subject to distinct market dynamics; currently, logs are valued at \$0.19/kg, while residues trade at a significantly lower price of \$0.11/kg. Given this vast scale and the sharp cost disparity between high-quality timber and lower-cost residues, static pricing data is insufficient for long-term strategic planning. Consequently, a robust price prediction model is required to accurately forecast market fluctuations and optimize feedstock procurement strategies.

The first model treats forest-residue biomass and log as interchangeable sources for products with wood fiber. 
This approach defines the supply curve using two key data points: the total available residues at a low price and the amount of residues and log supplies added at the current log price. 
The supply equation parameters are set to make the curve passes through both data points. 
However, this method often results in a supply curve that is more price-sensitive (\textit{elastic}) than what is typically observed in real markets, where pulpwood supply is usually less responsive to price changes. 

To address this discrepancy, the second approach uses a price elasticity adopted from previous research \cite{galik2009forest}.
This method assumes a constant, inelastic price response for pulpwood in the states of North Carolina, South Carolina, and Virginia with a fixed elasticity of 0.3 in a supply curve. 
It also assumes that all available forest-residue biomass is used up before any log is purchased for feedstock. 
The supply curve is constructed by following the same linear curve of the first model until the current log price and quantity. 
From then on, it maintains the elasticity of 0.3, making the slope steeper (SI~S7). 

\subsubsection{Green ammonia and hydrogen supply}\label{sec:methods-nh3}

Green ammonia is assigned a base price of \$635/t NH$_3$, with sensitivity analysis spanning \$270/t, representing a projected future price, to \$1,000/t, representing a current price. Hydrogen produced through ammonia cracking replaces grey hydrogen at a specified substitution fraction. The grey-hydrogen reference is priced at \$1,499/t H$_2$ \citep{katebah2022}.

\subsubsection{CCS cost}\label{sec:methods-ccscost}

Capture costs are calculated as described in Section~\ref{sec:methods-ccs}. T\&S costs are applied to the sequestered CO$_2$ mass at a base rate of \$10/t CO$_2$ , with a sensitivity range of \$4--45/t CO$_2$. The T\&S cost contribution varies linearly with its unit price.

\subsubsection{CO$_2$ abatement cost}
CO$_2$ abatement cost was used to evaluate the trade-off between hydrogen supply and CCS. Its definition is as follows: 
    \begin{equation}
    \text{Abatement cost} =
    \frac{\Delta\mathrm{MJSP}}
    {\Delta\mathrm{GHG}\times E_{saf}},
    \end{equation}
where \(E=31.66\)~MJ/L is the volumetric energy content of SAF. Because producers must at least meet the compliance threshold set by CORSIA, the reported value represents the additional cost of abatement beyond the reference configuration. Abatement costs values are undefined for configurations at or above the GHG intensity of compliance threshold. 


\subsection{Life-cycle assessment}\label{sec:methods-lca}

\subsubsection{System boundary and biogenic carbon}\label{sec:methods-boundary}

The functional unit is 1~MJ of SAF, and the system boundary covers the well-to-wake (WTWa) life cycle, comprising the well-to-pump (WTP) and pump-to-wake (PTWa) stages. The PTWa stage, representing fuel combustion, is adopted from the Fischer--Tropsch jet-fuel reference. Environmental impacts are expressed as GHG-100, based on 100-year global warming potential, in g~CO$_2$e/MJ SAF. Life-cycle inventories are derived from Greenhouse gases, Regulated Emissions, and Energy use in Technologies (GREET) model \cite{wang2023greet}, since it coincides with the standard of CORSIA. 

The total forest residue production emission includes forest maintenance activities, harvesting, collecting, bailing, and extracting, such as milling and chipping \cite{han2011well}.
In the meantime, ILUC emission of the feedstock is not considered since it has low risk in generating ILUC emissions, which applies similarly to biomass from wastes and microalgae \cite{ICAO2023CORSIA}.
This total emission of the production process is allocated to forest residue feedstock based on mass fractions, except for the emissions induced by fertilizer or agrochemical inputs \cite{han2011well}. 

Ammonia was assumed to be produced by renewable energy in Australia \cite{lee2022techno} and transported via conventional ship fueled by liquefied natural gas. Applying the technology of ammonia-fueled ships \cite{schwarzkopf2023future}, its benefit was negligible ($<$ 1.2 g CO$_2$e/MJ), so it was not included in the results of this paper.
Heat utilities such as mid-pressure and low-pressure steam were counted by heat consumption of fuel, which is natural gas from North America, without considering transportation. 
Heat exchange network in the process requires some make-up boiler feed water, and it is counted as process water resources. 

The counting of emissions from the downstream product chain imitates that of conventional synthetic paraffinic kerosene fuel. 
SAF produced from the plant is distributed via multiple routes following the assumption of the domestic jet distribution over the U.S. \cite{xie2011assessment}.
After reaching at pump stations, the PTWa emission profile is assumed to be the same as Fischer--Tropsch jet fuel from biomass \cite{elgowainy2012life}.
SAF compliance is assessed against the CORSIA requirement that eligible SAF achieve at least a 10\% reduction in WTWa GHG-100 intensity relative to conventional jet fuel.

Biogenic carbon is treated under the carbon-neutrality assumption. Biogenic combustion emissions are tracked separately as \(\text{CO$_{2, biogenic}$}\), distinct from fossil CO$_2$, such that biomass uptake and subsequent combustion emissions cancel. Biogenic CO$_2$ diverted to CCS is instead counted as a negative emission.

% \subsubsection{Allocation}\label{sec:methods-alloc}

Co-product emissions are allocated among the products for the central LCA estimate. This approach avoids strong dependence on displacement credits for exported electricity and fertilizer, whose results can vary substantially with the assumed displaced products and market conditions. 

% \subsubsection{SAF compliance criterion}\label{sec:methods-corsia}

\subsection{Sensitivity analysis}\label{sec:methods-sens}

% Three complementary analyses are conducted.

% First, cost-emission trade-off of hydrogen supply and CCS was evaluated.
% The main plant is fixed at the knee-point scale of 0.41~MMt/yr SAF. The green-ammonia hydrogen substitution fraction (0--1) and CCS capture fraction (0--1.0) are varied over a \(6\times5\) grid, yielding 30 configurations. Green NH$_3$ hydrogen shares of 0, 10, 25, 50, 75, and 100\% are crossed with CCS capture fractions of 0, 25, 50, 75, and 100\% (Table~\ref{tab:S9-4}). The non-dominated configurations are identified by simultaneously minimizing MJSP and well-to-wake GHG-100 intensity.

%%%% TO APPENDIX
% MJSP is calculated using the full TEA for every grid point. GHG intensity is evaluated on the allocation basis using the Korean electricity grid mix (Section~\ref{sec:methods-alloc}). It is calculated from the simulated CCS cases and a constant ammonia-related emissions adjustment:
%     \begin{equation}
%     \mathrm{GHG}(a,c) =
%     \mathrm{GHG}_{\mathrm{CCS}}(c)
%     -a\left(\mathrm{GHG}_{\mathrm{grey}}-\mathrm{GHG}_{\mathrm{NH}_3}\right),
%     \end{equation}
%     where \(a\) denotes the green-ammonia substitution fraction and \(c\) denotes the CCS capture fraction. This formulation assumes that the upstream emissions avoided by replacing grey H$_2$ are independent of, and additive to, the process CO$_2$ captured by CCS.

% Uncertainty bands for MJSP are constructed by varying the biomass price across the supply-curve range at this plant scale (107.3--116.1~\$/t), from the fixed forest-residue reference price to the supply-curve price (Section~\ref{sec:methods-feed}).

Sensitivity of process economics was compared by computing SHAP attribution of different parameters. Sensitivity analysis was applied only when the compliance threshold was satisfied based on LCA. For each configuration, 128 Latin hypercube samples are generated over the specified parameter ranges and evaluated using the TEA model. A gradient-boosting regression surrogate is fitted to the resulting samples, with an in-sample \(R^2>0.99\) required. TreeSHAP values are then calculated from the surrogate to attribute each sample's MJSP to individual drivers.

The varied parameters are grey-hydrogen price (\$1,000--4,000/t), green-ammonia price (\$270--1,000/t), CAPEX inflation (\(\pm30\%\)), CO$_2$ T\&S cost (\$4--45/t for CCS configurations), electricity price (\(\pm30\%\)), co-product prices (\(\pm20\%\)), and capacity factor (85\% or 91\%). Biomass price is additionally varied across the supply-curve range according to their production scale.

\begin{table}[htbp]
\centering
\caption{Varied parameters and corresponding ranges or values.}
\label{tab:varied_parameters}
\begin{tabular}{ll}
\hline
\textbf{Parameter} & \textbf{Range / Value} \\
\hline
Grey-H$_2$ price & \$1,000--4,000/t \\
Ammonia price & \$270--1,000/t \\
CAPEX & $\pm 30\%$ \\
CO$_2$ T\&S cost (CCS configurations) & \$4--45/t \\
Electricity price & $\pm 30\%$ \\
Co-product prices & $\pm 20\%$ \\
Capacity factor & 85\% or 91\% \\
Biomass price & Supply-curve range according to production scale \\
\hline
\end{tabular}
\end{table}
